\origpage{382}{N. Robertson, P.D. Seymour / Journal of Combinatorial Theory, Series B 90 (2004) 325--385}
\noindent similarly oriented, and so we may assume that $\Sigma_\chi$ is not orientable. Since $\Sigma_\chi$ is connected, there is, for each internal $e \in E(\Gamma_0)$ with $|\tilde{e}| = 3$, a $\Phi_\chi$-preserving homeomorphism $\alpha_e \colon \Sigma_\chi \to \Sigma_\chi$ which maps $\beta(e)$ onto itself with reversed orientation, and fixes $bd(\Sigma)$ and every $\beta(f)$ ($f \in E(\Gamma_0)\setminus \{e\}$) pointwise. By an appropriate combination of the $\alpha_e$'s we may correct every $e \in E(\Gamma_0)$ with $|\tilde{e}| = 3$ for which $\alpha(\gamma_\Gamma(\beta(e)))$ and $\gamma_{\Gamma_0}(e)$ give opposite orientation of $\beta(e)$. This proves (2).

\hspace*{12pt}(3) \textit{For each $v \in V(\Gamma_0)$ there is a tree $T_v$ in $\Sigma_\chi$ such that}

\hspace*{12pt}$\bullet$ \textit{for distinct $v, v' \in V(\Gamma_0)$, $U(T_v) \cap U(T_{v'}) = \emptyset$,}

\hspace*{12pt}$\bullet$ \textit{for each $v \in V(\Gamma_0)$, $e \in E(\Gamma_0)$ and $1 \le i \le |\tilde{e}|$, $U(T_v)$ contains the $i$th term of $\gamma_\Gamma(\beta(e))$ if and only if $v$ is the $i$th term of $\gamma_{\Gamma_0}(e)$, and}

\hspace*{12pt}$\bullet$ \textit{for each $v \in V(\Gamma_0)$ and $e \in E(\Gamma_0)$, $U(T_v) \cap \beta(e) = \emptyset$.}

\textbf{Subproof.} Let $S_v$ be a tree in $\Sigma_\chi$ for each $v \in V(\Gamma_0)$, such that

\hspace*{12pt}$\bullet$ for distinct $v, v' \in V(\Gamma_0)$, $U(S_v) \cap U(S_{v'}) = \emptyset$,

\hspace*{12pt}$\bullet$ for each $v \in V(\Gamma_0), e \in E(\Gamma_0)$ and $1 \le i \le |\tilde{e}|$, $U(S_v)$ contains the $i$th term of $\alpha(\gamma_\Gamma(\beta(e)))$ if and only if $v$ is the $i$th term of $\gamma_{\Gamma_0}(e)$, and

\hspace*{12pt}$\bullet$ for each $v \in V(\Gamma_0)$ and $e \in E(\Gamma_0)$, $U(S_v) \cap \alpha(\beta(e)) = \emptyset$.

These $S_v$'s clearly exist (let each $S_v$ be a star centered at $v$, with edges entering all those $e \in E(\Gamma_0)$ with $v \in \tilde{e}$ except short sides). Let $T_v = \alpha^{-1}(S_v)$ for each $v \in V(\Gamma_0)$; then (3) is satisfied. This proves (3).

\fix{The rest of the proof uses (2) and (3) through the topological feasibility condition of Theorem 3.2 of [7], which asks for closed discs in place of trees. We give the details of (2) and (3) in that form. Theorem 3.2 of [7] is applied with the following closed discs in $\hat{\Sigma}_\chi$: for each cuff $\Theta$, with pole $r^*$ say, a disc $\Delta_\Theta$ in the closure of the disc pasted onto $\Theta$, with $\Delta_\Theta \cap \Sigma_\chi = N(r^*)$; and for each internal $e \in E(\Gamma_0)$ a disc $\Delta_e$ that meets $sk(\Gamma)$ exactly in the ends of $\beta(e)$, all in $bd(\Delta_e)$, and that lies in the closure of $\beta(e)$ if $|\tilde{e}| = 3$. These discs may be chosen mutually disjoint, since by (1) the edges $\beta(e)$ are far from each other and from the poles. For $v \in V(\Gamma_0)$ let $Z_v$ be the set of all $i$th terms of $\gamma_\Gamma(\beta(e))$ such that $v$ is the $i$th term of $\gamma_{\Gamma_0}(e)$, over all $e \in E(\Gamma_0)$ and $i$. We show that there are mutually disjoint closed discs $D_v$ $(v \in V(\Gamma_0))$ such that $D_v \cap \bigcup \Delta = Z_v$ for each $v$, where $\bigcup \Delta$ denotes the union of all the discs $\Delta_\Theta$ and $\Delta_e$.

For each of these discs $\Delta$ we choose a closed disc $P$, its \textit{pad}, in which $\Gamma_0$ plays the role of $\Gamma$: for $\Delta_\Theta$ the closure of the disc pasted onto $\Theta$; for $\Delta_e$ with $|\tilde{e}| = 3$ the closure $\bar{e}$; and for $\Delta_e$ with $|\tilde{e}| = 2$ a thin closed disc including $\bar{e}$, with the ends of $e$ in its boundary. The pads may be chosen to meet only in $V(\Gamma_0)$, with every vertex of $\Gamma_0$ in a pad in its boundary. Each point of $Z_v$ in $bd(\Delta)$ corresponds to the vertex $v$ in $bd(P)$, and these points and vertices appear in the same cyclic order around $bd(\Delta)$ and $bd(P)$ (several points may correspond to one vertex), for suitable directions of the two circles; and if $\Sigma_\chi$ is orientable, these directions can be chosen to agree with one orientation of $\hat{\Sigma}_\chi$ for all the discs at once. For the cuff discs this holds because by (1) the edges $\beta(f_1), \dots, \beta(f_k)$ border each long side in the order of $f_1, \dots, f_k$, and $\Delta_\Theta$ lies in the pad of $\Theta$; for $|\tilde{e}| = 3$ it holds because $(\Gamma_0, \phi_0)$ and $(\Gamma, \phi)$ are similarly oriented and $\Delta_e \subseteq \overline{\beta(e)}$ (and if $\Sigma_\chi$ is not orientable there is nothing to prove, since three points of a circle are in the given cyclic order for one of the two directions); and for $|\tilde{e}| = 2$ there is nothing to prove.

Now choose small mutually disjoint closed discs $E_v$ $(v \in V(\Gamma_0))$, with $v$ in the interior of $E_v$, each meeting each pad only near $v$. In each pad $P$ choose a closed disc $M_P$ that meets $\bigcup_v E_v$ in finitely many points, its \textit{peaks}, one for each point of $Z_v$ in the corresponding disc $\Delta$, lying in $E_v$ near $v$, and with the peaks in the same cyclic order around $bd(M_P)$ as these points around $bd(\Delta)$. (In a chart of $P$ as the unit disc, take $M_P = \{z \colon |z| \le \rho(z/|z|)\}$, where $\rho < 1$ except near the directions of the vertices of $\Gamma_0$, where $\rho$ rises just enough to reach the discs $E_v$.) The discs $M_P$ are mutually disjoint, since the pads meet only in $V(\Gamma_0)$ and the $M_P$ miss $V(\Gamma_0)$. Next there is a homeomorphism $\alpha'$ of $\hat{\Sigma}_\chi$ with $\alpha'(M_P) = \Delta$ for every pad $P$ of a disc $\Delta$, mapping each peak to its point of $Z_v$. To see this, take first a homeomorphism mapping each $M_P$ onto its disc $\Delta$, which exists since both are families of the same number of disjoint closed discs in a connected surface; if $\Sigma_\chi$ is orientable we choose it to preserve the orientation. On each boundary circle the remaining error is a homeomorphism of the circle. If it preserves the direction of the circle, a twist in a thin annulus around the disc corrects it, without moving the other discs. If it reverses the direction, then $\Sigma_\chi$ is not orientable, since otherwise the directions of the cyclic orders agree with the orientation; and then there is a homeomorphism of $\hat{\Sigma}_\chi$ that maps this disc onto itself with the direction of its boundary reversed and fixes the other discs, after which a twist corrects the error. Finally let $D_v = \alpha'(E_v)$ for each $v \in V(\Gamma_0)$. These discs are mutually disjoint, and $D_v \cap \bigcup \Delta = \alpha'(E_v \cap \bigcup_P M_P) = Z_v$, as required. In the terms of (2) and (3), the map $\alpha'^{-1}$ plays the role of $\alpha$, the disc $M_P$ that of $\alpha(\beta(e))$, and $E_v$ that of $S_v$.}\fixnote{fix:15.2-Sv-proof-added}{The paper says that the $S_v$ clearly exist, and that most of (2) is easy to arrange. The text above gives (2) and (3) in the form of closed discs, which is the form used by Theorem 3.2 of [7]. This is the argument used in the formalization.}

\hspace*{12pt}(4) \textit{Let $r^*$ be a pole, and let $N(r^*) = \bigcup (\tilde{\beta}(e) \colon e \in E(\Gamma_0), e \subseteq \bar{r}\setminus r)$. There is no circuit $F$ of $\Gamma^*$ with $\frac{1}{2}|E(F)| < rep(\mathscr{T})$ and with $r^* \in ins(F)$ and $\frac{1}{2}|E(F)| < |ins(F) \cap N(r^*)|$.}

\textbf{Subproof.} Suppose that $F$ is such a circuit. Now $|N(r^*)| \le 2|V(\Gamma_0)|$, and so
\begin{equation*}
\frac{1}{2}|E(F)| < 2|V(\Gamma_0)| \le \frac{1}{2}n - 2.
\end{equation*}
Suppose first that $r^* \in ins(F)\setminus U(F)$. Since $\frac{1}{2}|E(F)| \le n$, and $\bar{r} \subseteq ins(F)$, it follows that $\bar{r}$ includes no long side of $\Phi_\chi$, because the second outcome of the theorem does not hold. Hence $N(r^*) = V(\Phi_\chi) \cap (\bar{r}\setminus r)$, and so $ins(F)$ contains more than $\frac{1}{2}|E(F)|$ vertices of $\Phi_\chi$, and the third outcome of the theorem holds, a contradiction.

Thus $r^* \notin ins(F)\setminus U(F)$, and so $r^* \in V(F)$. Hence $F$ is a bite at $r^*$. We may assume that $ins(F) \cap (\bar{r}\setminus r)$ is a line, for otherwise there is a circuit $F'$ with $|E(F')| < E(F)$ and $ins(F) \subseteq ins(F')$. Suppose that $\beta(e) \subseteq ins(F)$ for some $e \in E(\Gamma_0)$ bordering $\bar{r}\setminus r$ which is not a short side. Let $S$ be the long side with $e \subseteq S$, and let $s_1, s_2$ be the short sides with common ends with $S$. (Possibly $s_1 = s_2$.) Let $v_1, v_2$ be the corresponding ends of $S$. If $v_1 \in ins(F)$ then $l(\beta(e), s_1) \le \frac{1}{2}|E(F)| + 2 \le \frac{1}{2}n$, a contradiction to the third assertion of (1). Thus $v_1, v_2 \notin ins(F)$. Since $ins(F) \cap (\bar{r}\setminus r)$ is a line, it follows that $ins(F) \cap (\bar{r}\setminus r) \subseteq S$. By (1) again, $\beta(f) \nsubseteq ins(F)$ for any $f \in E(\Gamma_0)\setminus \{e\}$ which borders $S$, and so $|ins(F) \cap N(r^*)| = 2$. Hence $\frac{1}{2}|E(F)| \le 1$, which is impossible.
